\section{Non-rigid Hall objects and their theta straightening}
\label{sec:nonrigid}

The subquotient theorem is stronger than a statement about generators.  It
assigns a canonical quantum cluster element to every Hall basis object,
including objects with self-extensions.  This section gives the explicit
recursive formula and proves its independence of choices.

\subsection{The Hall-corrected straightening operator}

Let $M$ be a reduced object and write its multidiagonal as
\[
 D(M)=\{\gamma_1,\ldots,\gamma_r\}.
\]
If $D(M)$ is non-crossing, set
\begin{equation}
 \St_{\mathrm H}(M)=\eps_M.
 \label{eq:straightening-rigid-base}
\end{equation}
Suppose $D(M)$ has a crossing.  Choose crossing summands
$M_\gamma,M_\delta$ and write
\[
 M=M_\gamma\oplus M_\delta\oplus R.
\]
Expand the actual Hall product in the fixed order:
\begin{equation}
 (\eps_\gamma\star_S\eps_\delta)\star_S\eps_R
 =a_M\eps_M+
  \sum_{E\prec M}a_E\eps_E,
 \qquad a_M\ne0.
 \label{eq:Hall-product-containing-M}
\end{equation}
The coefficient $a_M$ is a product of split Hall coefficients and is
invertible in $\mathbb Q(\omega)$.  Replace the crossing pair by the skein
right-hand side and define recursively
\begin{equation}
\boxed{
\begin{aligned}
 \St_{\mathrm H}(M)=a_M^{-1}\Big(&
 \big(
   \omega^{a(\gamma,\delta)}
      \eps_{\gamma\#_0\delta}
  +\omega^{b(\gamma,\delta)}
      \eps_{\gamma\#_\infty\delta}
 \big)\star_S\eps_R\\
 &-\sum_{E\prec M}a_E\St_{\mathrm H}(E)
 \Big).
\end{aligned}}
\label{eq:Hall-corrected-straightening}
\end{equation}
Every recursive argument has smaller crossing number by
\cref{thm:extension-smoothing}.

\begin{theorem}[Canonical Hall straightening]
\label{thm:canonical-Hall-straightening}
The recursion \eqref{eq:Hall-corrected-straightening} terminates and is
independent of:
\begin{enumerate}[label=(\roman*),leftmargin=2.3em]
\item the chosen crossing pair;
\item the chosen parenthesization of the Hall product;
\item the chosen sequence of subsequent crossings.
\end{enumerate}
It yields a finite expression
\begin{equation}
 \St_{\mathrm H}(M)
 =\sum_{D\in\mathfrak N(P)}c_{M,D}(\omega)\eps_D,
 \qquad c_{M,D}(\omega)\in\mathbb Q(\omega),
 \label{eq:Hall-rigid-expansion}
\end{equation}
which is the unique normal-form representative of $\eps_M$ modulo $\Jsk$.
\end{theorem}

\begin{proof}
Termination is \cref{lem:termination}.  Formula
\eqref{eq:Hall-corrected-straightening} is exactly the oriented reduction
obtained by multiplying a Hall--skein defect by the remaining factors and
solving for its invertible leading term.  Choice independence and uniqueness
are therefore the confluence statement of
\cref{thm:Hall-skein-diamond}.  The irreducible terms are the non-crossing
basis classes of \cref{cor:rigid-normal-basis}.
\end{proof}

The recursion is effective whenever the finite-field Hall coefficients for
the relevant presentation spaces can be computed.  It never asks for a
quantum Euler characteristic of a singular quiver Grassmannian: all quantum
coefficients are obtained from exact Hall numbers and the fixed local skein
relation.

\subsection{Theta interpretation}

Compose the normal form with the global quotient map.  For every reduced
object $M$, define
\begin{equation}
 \qchar_M^{\rm sk}
 :=\Phi_{\rm sk}(\eps_M).
 \label{eq:skein-quantum-character}
\end{equation}
This is a well-defined quantum cluster element for rigid and non-rigid $M$.
It should be distinguished from an a priori quiver-Grassmannian quantum
cluster character, which is not known in this generality.

\begin{theorem}[Non-rigid theta expansion]
\label{thm:nonrigid-theta-expansion}
For every reduced Hall basis object $M$,
\begin{equation}
 \boxed{
 \qchar_M^{\rm sk}
 =\sum_{D\in\mathfrak N(P)}
    c_{M,D}(\omega)\Theta_{g(D)},
 }
 \label{eq:nonrigid-theta-expansion}
\end{equation}
where the coefficients are those of
\eqref{eq:Hall-rigid-expansion}.  The expansion is finite and unique.
Moreover:
\begin{enumerate}[label=(\roman*),leftmargin=2.3em]
\item if $M$ is rigid, the sum has one term and
      $\qchar_M^{\rm sk}=\Theta_{g(M)}$;
\item if $M$ is non-rigid, the expansion has support only on rigid direct sums
      with the same frozen/weight degree;
\item no positivity or bar-invariance of the individual coefficients
      $c_{M,D}(\omega)$ is asserted.
\end{enumerate}
\end{theorem}

\begin{proof}
Apply $\Phi_{\rm sk}$ to
\eqref{eq:Hall-rigid-expansion}.  By
\cref{thm:Hall-skein-subquotient}, the image of a non-crossing class $\eps_D$
is $\Theta_{g(D)}$.  Uniqueness follows because the non-crossing skein basis is
the theta basis in the polygon case.
\end{proof}

In finite type, the cluster fan is complete and every theta label belongs to a
cluster cone.  Thus the right-hand side of
\eqref{eq:nonrigid-theta-expansion} is also the unique expansion in normalized
quantum cluster monomials.

\subsection{Incompatible rigid products}

Let $X,Y$ be rigid but incompatible.  Their individual images are theta basis
elements, but their exact Hall product contains non-rigid middle terms.  Write
\[
 \eps_X\star_S\eps_Y
 =\sum_E h_{X,Y}^E\eps_E.
\]
Substituting \eqref{eq:nonrigid-theta-expansion} gives the explicit identity
\begin{equation}
 \Theta_{g(X)}\Theta_{g(Y)}
 =\sum_{E,D}
   h_{X,Y}^E(\omega)c_{E,D}(\omega)\Theta_{g(D)}.
 \label{eq:incompatible-theta-expansion}
\end{equation}
The total coefficient of $\Theta_{g(D)}$ is therefore
\begin{equation}
 m_{X,Y}^{D}(\omega)
 =\sum_E h_{X,Y}^E(\omega)c_{E,D}(\omega).
 \label{eq:theta-structure-from-Hall}
\end{equation}
These are exactly the theta multiplication structure constants, because the
left-hand side is the product in the quantum cluster algebra.

\begin{corollary}[Quantum meaning of every middle term]
\label{cor:middle-term-meaning}
Every non-rigid middle term $E$ appearing in an incompatible Hall product has
a canonical, finite quantum/theta value $\qchar_E^{\rm sk}$.  The weighted sum
of these values with the exact Hall coefficients is the theta product.  Hence
no Hall stratum is discarded; the quotient identifies the linear combinations
forced by the skein relations.
\end{corollary}

\subsection{The one-crossing case}

Suppose $\gamma,\delta$ cross once.  The two $2$-Calabi--Yau extension
directions have middle terms the two smoothings.  The exact morphism Hall
product sees the split object and the exact extensions in its chosen
direction.  Write
\[
 \eps_\gamma\star_S\eps_\delta
 =h^{\rm sp}\eps_{\gamma\oplus\delta}
  +\sum_jh_j\eps_{E_j}.
\]
The quotient relation gives the explicit non-rigid value
\begin{equation}
 \qchar_{\gamma\oplus\delta}^{\rm sk}
 =(h^{\rm sp})^{-1}\left(
   \omega^a\Theta_{g(\gamma\#_0\delta)}
  +\omega^b\Theta_{g(\gamma\#_\infty\delta)}
  -\sum_jh_j\qchar_{E_j}^{\rm sk}
 \right).
 \label{eq:one-crossing-nonrigid}
\end{equation}
Thus the missing opposite smoothing is not claimed to be an exact extension in
the same Hall direction.  It enters through the quotient relation dictated by
the $2$-Calabi--Yau/skein multiplication law.

\subsection{Relation with cluster multiplication formulas}

For a pair of objects $L,M$ with one-dimensional extension space in a
$2$-Calabi--Yau category, the classical cluster character satisfies
\[
 CC_T(L)CC_T(M)=CC_T(E)+CC_T(E'),
\]
where $E,E'$ are the two exchange middle terms.  Caldero--Keller prove this in
finite type, and Palu proves the general constructible $2$-Calabi--Yau formula
\cite{CalderoKeller,PaluMultiplication}.  Chen--Xiao--Xu and
Xiao--Xu--Yang give weighted quantum and motivic refinements for arbitrary
pairs \cite{ChenXiaoXu,XiaoXuYang}.

Our quotient may be viewed as a strict algebraic packaging of these formulas
in polygon type.  Rather than assigning a character only after selecting a
pair and summing extension varieties, it imposes all local two-direction
multiplication identities as one two-sided ideal in the exact morphism Hall
algebra.  The Diamond theorem proves that the identities are globally
consistent and that their quotient has the expected theta basis.

\subsection{Algorithmic form}

\begin{algorithm}[Hall-to-theta straightening]
\label{alg:Hall-to-theta}
Given a reduced Hall object $M$:
\begin{enumerate}[label=\arabic*.,leftmargin=2.3em]
\item Compute its Krull--Schmidt multidiagonal $D(M)$.
\item If $D(M)$ is non-crossing, output $\Theta_{g(M)}$.
\item Choose a crossing pair $\gamma,\delta$ and compute the exact Hall
      expansion \eqref{eq:Hall-product-containing-M}.
\item Apply \eqref{eq:Hall-corrected-straightening}.
\item Recurse on all lower-crossing middle terms and collect the theta basis
      coefficients.
\end{enumerate}
The number of recursive levels is at most $\crs(M)$.
\end{algorithm}

The accompanying code implements the geometric straightening and confluence
part.  Computing the exact coefficients $a_E$ from an arbitrary bound quiver
is a separate representation-theoretic task and is not silently replaced by
polygon combinatorics.
